\documentclass{article}
\usepackage{amsmath,amsthm}
\ifdefined\ReviewMode
  \edef\LoadReviewPackage{\noexpand\usepackage[\ReviewMode]{advisorreview}}%
  \LoadReviewPackage
\else
  \newcommand{\pdfrevnote}[1]{}
  \newcommand{\pdfrevmark}[2]{#1}
  \newcommand{\pdfrevmath}[2]{#1}
  \newcommand{\pdfrevformula}[2]{#1}
  \newcommand{\pdfrevblock}[2]{#1}
\fi
\newtheorem{theorem}{Theorem}
\pagestyle{empty}
\newcommand{\measuremath}[2]{%
  \setbox0=\hbox{\(\displaystyle #2\)}%
  \typeout{REVIEW-DIM:#1:\the\wd0:\the\ht0:\the\dp0}%
}
\begin{document}
\measuremath{relation}{a\pdfrevformula{=}{Measured relation}b}
\measuremath{binary}{a\pdfrevformula{+}{Measured operator}b}
\measuremath{unary}{(\pdfrevformula{+}{Measured unary plus}b)}
\measuremath{limits}{\pdfrevformula{\sum}{Measured sum}\limits_{k=1}^n k}
\measuremath{scripts}{\pdfrevformula{x}{Measured subscript}_k+\pdfrevformula{x}{Measured superscript}^2}
\noindent Formula fragments, with ordinary TeX spacing:
\[
  A_e\left(\bar e(t)+\pdfrevformula{\bar x}{State correction}\right)\,dt+D_e\,d\omega
\]
\[
  e(t)+\int_0^r A_\zeta\left(\bar e(\pdfrevformula{s}{Integration variable})+\bar x\right)\,ds
\]
\[
  W(\mathcal G_1(q))\pdfrevformula{=}{Equality}\lambda W(q),\qquad
  \pdfrevformula{D^+}{Dini derivative}W\leq 2\sqrt{W}\,\widetilde H_1(\bar e)
\]
\[
  a\pdfrevformula{+}{Binary operator}b\pdfrevformula{\leq}{Relation}c,
  \qquad \pdfrevformula{\sin}{Function}(x)+\pdfrevformula{\sum}{Sum}_{k=1}^n k
\]
\[
  \pdfrevformula{x}{External subscript}_k+
  \pdfrevformula{x}{External superscript}^2+
  x_{\pdfrevformula{i+1}{Script style}}+
  x^{y_{\pdfrevformula{j+1}{Nested script style}}}
\]
\[
  (\pdfrevformula{+}{Unary plus}b),\quad
  \pdfrevformula{\sum}{Explicit limits}\limits_{k=1}^n k,\quad
  \pdfrevformula{\int_0^1}{Integral} x\,dx,\quad
  a\pdfrevmark{=}{Math-mode text macro}b
\]
\begin{equation*}
  \check y_k(t_k^+)=y(t_k)+h_y\left(k,\pdfrevformula{\check y_{k-1}(t_k-y(t_k))}{Corrected symbol}\right)
\end{equation*}
\begin{align*}
  a&\pdfrevformula{=}{Alignment relation}b+\pdfrevformula{c}{Alignment variable},\\
  x&\leq \pdfrevformula{\frac{1}{1+t}}{Fraction}.
\end{align*}
\noindent Inline: \pdfrevmath{\(\alpha_k>0\)}{Positivity} and
\pdfrevmark{short revised text}{Wording}, punctuation stays here.

\noindent Sentence before. \pdfrevblock{This sentence stays in the same
paragraph and contains \(x_k\), a citation-like reference, and enough words
to exercise line wrapping at the normal text width.}{Sentence revision}
Sentence after.

\begin{theorem}
\pdfrevblock{If \(x=0\), then \(x^2=0\).}{Theorem content}
\end{theorem}
\noindent Text after the theorem.\pdfrevnote{Standalone note}
Sentence spacing after a note stays unchanged.
\end{document}
